\documentclass{beamer}
\usepackage[british,spanish]{babel}
\usepackage[utf8]{inputenc}
\usepackage[british,spanish]{babel}
%usepackage[utf8]{inputenx}

\usepackage{hyperref}
%\hypersetup{colorlinks=false,linkbordercolor=red,linkcolor=green,pdfborderstyle={/S/U/W 1}}

\usepackage{multirow}

\usepackage{textcomp}

\usepackage{listings}
\lstloadlanguages{Ruby}
\usepackage{cancel}

\usepackage{adjustbox}
\usepackage{lstcustom}

\usepackage{amsmath}

\usepackage{color}
\definecolor{light-gray}{gray}{0.80}
\definecolor{lstbackgroundshellcolor}{named}{light-gray}

\usepackage{tikz}
\newcommand*\circled[1]{\tikz[baseline=(char.base)]{
            \node[shape=circle,draw,inner sep=2pt] (char) {#1};}}

\usepackage[normalem]{ulem}

%\usepackage[acronym,xindy,toc]{glossaries}

\usepackage[acronym,xindy,toc]{glossaries}
\makeglossaries
%\usepackage[xindy]{imakeidx}
%\makeindex


\newcommand{\comment}[2]{#2}

\newcommand{\commandinline}[1]{\lstinline[basicstyle=\small\lstfontfamily]{#1}}
\newcommand{\outputcommand}[1]{\color{darkgreen}{#1}}

\graphicspath{ {./images/} }

\title{Amazon SimpleDB: A Cloud Database}
%\subtitle[short subtitle]{long subtitle}
\author[C. Cuenca, F. Quintana]{Carmelo Cuenca-Hernández and Francisca Quintana-Domínguez}
%\institute{Escuela Universitaria de Informática}
%\date[04/2013]{Abril - 2013}
\date{}
\titlegraphic{\includegraphics[width=0.5 \textwidth]{./images/logo_ulpgc_version_horizontal_rgb.eps}}


\pgfdeclareimage[width=2.0\baselineskip]{ulpgc-logo}{images/logosimbolo_secundario_version_vertical}
\setbeamertemplate{footline}{\raisebox{-2ex}{\pgfuseimage{ulpgc-logo}}
  \usebeamerfont{date in head/foot}\insertshortdate{}\hfill
  \usebeamertemplate{navigation symbols}\hfill
  \insertframenumber{}/\inserttotalframenumber}
\setbeamertemplate{sidebar right}{}


\usetheme{Antibes}
%\usetheme{Berlin}

%\usetheme{Warsaw}
%\usecolortheme{albatross}

\selectlanguage{british}



\begin{document}

%\includegraphics[width= 1.0 \textwidth]{logos3.eps}
\begin{frame}
	\titlepage
\end{frame}


\section*{Outline}
\begin{frame}[fragile, allowframebreaks]
  \frametitle{Outline}
  %\tableofcontents%[part=1,pausesections]
  \tableofcontents[currentsection,currentsubsection, sectionstyle=show] 
  %\tableofcontents[currentsection,sectionstyle=show,hideothersubsections]
\end{frame}

%%%%%%%%%%%%%%%%%%%%%%%%%%%%%%%%%%%%%%%%%%%%%%%%%%%%%%%%%%%%%%%%%%%%%%%%%%%%%%
%\newacronym{<label>}{<abbrv>}{<full>}
%\glsreset{<label>}
%\glsresetall
%\acrlong{<label>}
%\acrfull{<label>}
%\acrshort{<label>}
%\input{../glossary}

\newacronym{acl}{ACL}{Access Control List}
\newacronym{api}{API}{Application Programming Interface}
\newacronym{aws}{AWS}{Amazon Web Services}
\newacronym{cli}{CLI}{Command Line Interface}
\newacronym{css}{CSS}{cascading style sheets}
\newacronym{ebs}{EBS}{Elastic Block StoAWS Identity an Access Managementrage}
\newacronym{ec2}{EC2}{Amazon Elastic Compute Cloud}
\newacronym{elb}{ELB}{Elastic Load Balancing}
\newacronym{iam}{IAM}{Identity Access Management}
\newacronym{ror}{RoR}{{\href{http://rubyonrails.org/}{Ruby on Rails}}}
\newacronym{rds}{RDS}{Relational Database Service}
\newacronym{rvm}{RVM}{{\href{https://rvm.io/}{Ruby Version Manager}}}
\newacronym{s3}{S3}{Simple Storage Service}
\newacronym{sqs}{SQS}{Amazon Simple Queue Service}
%%%%%%%%%%%%%%%%%%%%%%%%%%%%%%%%%%%%%%%%%%%%%%%%%%%%%%%%%%%%%%%%%%%%%%%%%%%%%%
\section{Amazon SimpleDB}
\begin{frame}[fragile]
\frametitle{Amazon SimpleDB}
\begin{itemize}
 \item Amazon SimpleDB is a cloud-based database
 \item SimpleDB excels at storing semi-structured data where the items (rows) are similar but not necessarily identical to each other
 \item All data stored in Amazon SimpleDB is automatically indexed
 \item The Amazon SimpleDB model lacks support for joins across domains (roughly
equivalent to relational tables); instead, you can store data in non-normalized form
for more efficient access
\end{itemize}
\end{frame}
%%%%%%%%%%%%%%%%%%%%%%%%%%%%%%%%%%%%%%%%%%%%%%%%%%%%%%%%%%%%%%%%%%%%%%%%%%%%%%
\begin{frame}[fragile]
\frametitle{Amazon SimpleDB Concepts}
\begin{itemize}
 \item An Amazon SimpleDB \emph{domain} is roughly analogous to a table in a relational database.  Each domain can store up to \texttt{10GB} of data
 \item SimpleDB excels at storing semi-structured data where the items (rows) are similar but not necessarily identical to each other
 \item Each item in a SimpleDB domain has a name (unique to the domain) and up to \texttt{256} attributes (name-value pairs).
       Item names, attribute names, and attribute values can each be up to \texttt{1024} bytes long
 \item Item attributes can be multivalued 
 \item All attribute values are treated as strings (\alert{numerical values?})
\end{itemize}
\end{frame}

%%%%%%%%%%%%%%%%%%%%%%%%%%%%%%%%%%%%%%%%%%%%%%%%%%%%%%%%%%%%%%%%%%%%%%%%%%%%%%
%%%%%%%%%%%%%%%%%%%%%%%%%%%%%%%%%%%%%%%%%%%%%%%%%%%%%%%%%%%%%%%%%%%%%%%%%%%%%%
\begin{frame}[fragile]
\frametitle{\texttt{Select}}
\begin{tabular}{|l|l|l|l|l|l|l|} \hline
item\_name & first\_name & last\_name & age & sex & middle & state \\ \hline
Rec1 & Tom & Basic & 12 & M & & \\ \hline
Rec2 & Nancy & Hacker & 15 & F & & MD \\ \hline 
Rec3 & Joan & Hughes & 44 & F & Fraft & NY, NJ\\ \hline 
\end{tabular} 

\begin{itemize}
 \item \texttt{select * from People where FirstName="Tom"} (Rec1)
 \item \texttt{select FirstName,Middle,LastName from People where LastName >= "H"} (Rec2 and Rec3)
 \end{itemize}
\end{frame}
%%%%%%%%%%%%%%%%%%%%%%%%%%%%%%%%%%%%%%%%%%%%%%%%%%%%%%%%%%%%%%%%%%%%%%%%%%%%%%
%%%%%%%%%%%%%%%%%%%%%%%%%%%%%%%%%%%%%%%%%%%%%%%%%%%%%%%%%%%%%%%%%%%%%%%%%%%%%%
\section{Amazon SimpleDB Pricing}
\begin{frame}[fragile]
\frametitle{Amazon SimpleDB Pricing}

\begin{itemize}
 \item Data transfer. Data transferred into SimpleDB is charged at a rate of \$0.10 per gigabyte
 \item Data storage. You pay \$0. 25 ( a quarter) per gigabyte per month to store data in Simple DB
 \item Machine utilization. You pay \$0.14 per hour for the machine time used to process each SimpleDB
request
 \end{itemize}
\end{frame}

%%%%%%%%%%%%%%%%%%%%%%%%%%%%%%%%%%%%%%%%%%%%%%%%%%%%%%%%%%%%%%%%%%%%%%%%%%%%%%
\section{Amazon SimpleDB Programming Model}
\begin{frame}[fragile]
\frametitle{Amazon SimpleDB Programming Model}
\begin{itemize}
 \item The Amazon SimpleDB programming model consists of just nine calls:
 \begin{itemize}
 \item At the domain level: \texttt{CreateDomain}, \texttt{DeleteDomain}, \texttt{ListDomains}, \texttt{DomainMetadata}
 \item At the item level
 \begin{itemize}
 \item \texttt{PutAttributes} creates new items and adds or replaces (your choice) additional attributes to existing items
 \item \texttt{BatchPutAttributes} is an extended version of this call that handles multiple items at once
 \item \texttt{DeleteAttributes} removes attributes from an item
 \item \texttt{GetAttributes} retrieve specified attributes
 \item \texttt{Select} issues a SQL query
 \end{itemize}
\end{itemize}

\end{itemize}
\end{frame}

%%%%%%%%%%%%%%%%%%%%%%%%%%%%%%%%%%%%%%%%%%%%%%%%%%%%%%%%%%%%%%%%%%%%%%%%%%%%%%
%%%%%%%%%%%%%%%%%%%%%%%%%%%%%%%%%%%%%%%%%%%%%%%%%%%%%%%%%%%%%%%%%%%%%%%%%%%%%%
\begin{frame}[fragile, allowframebreaks]
\frametitle{Create a Domain}
\begin{itemize}
\item Here’ s a script that creates all the SimpleDB domains needed for the examples
\lstset{language=Ruby, style=eclipse}
\begin{lstlisting}[escapechar=!]
#
# aws_app/simpledb/book.inc.rb
#
# This file is required by some *.rb files
# It defines default constant value

BOOK_FILE_DOMAIN = 'files'
BOOK_FEED_DOMAIN = 'feeds'
BOOK_FEED_ITEM_DOMAIN = 'feed_items'
\end{lstlisting}
\lstset{language=Ruby, style=eclipse}
\begin{lstlisting}[escapechar=!]
#
# aws_app/simpledb/create_domain.rb
#
# Usage: ruby simpledb/create_domain.rb
#

require File.expand_path("#{File.dirname(__FILE__)}/../config")
require File.expand_path("#{File.dirname(__FILE__)}/book.inc")

#get an instance of the SDB interface using the default configuration
sdb = AWS::SimpleDB.new(region: 'us-east-1')

[BOOK_FILE_DOMAIN, BOOK_FEED_DOMAIN, BOOK_FEED_ITEM_DOMAIN].each do |domain_name|
  begin
    domain = sdb.domains.create(domain_name)
      puts "Domain #{domain.name} created"
  rescue Exception => e
    puts e.message
  end
end
\end{lstlisting}

\item Modify the above script to accept new domain names on the command line (Exercise)
\item Install SDBTool plugin in Firefox from http://code.google.com/p/sdbtool/ and use it (Exercise)
\end{itemize}

\end{frame}

%%%%%%%%%%%%%%%%%%%%%%%%%%%%%%%%%%%%%%%%%%%%%%%%%%%%%%%%%%%%%%%%%%%%%%%%%%%%%%
\begin{frame}[fragile]
\frametitle{List Domains}

\lstset{language=Ruby, style=eclipse}
\begin{lstlisting}[escapechar=!]
#
# aws_app/simpledb/list_domains.rb
#
# Usage: ruby simpledb/lists_domains.rb
#

require File.expand_path(File.dirname(__FILE__) + '/../config')
require File.expand_path(File.dirname(__FILE__) + '/book.inc')

# get an instance of the SDB interface using the default configuration
sdb = AWS::SimpleDB.new region: 'us-east-1'

sdb.domains.each do |domain|
  puts domain.name
end
\end{lstlisting}

\end{frame}

%%%%%%%%%%%%%%%%%%%%%%%%%%%%%%%%%%%%%%%%%%%%%%%%%%%%%%%%%%%%%%%%%%%%%%%%%%%%%%
\begin{frame}[fragile, allowframebreaks]
\frametitle{Storing Data}
\begin{itemize}
\item The next step is to store some data in a domain. This is done using the
\texttt{put} method. This method can be used to create new items or to add
additional attributes to existing items

\lstset{language=Ruby, style=eclipse}
\begin{lstlisting}
#
# aws_app/simpledb/insert_items.rb
#
# Usage: ruby simpledb/insert_items.rb
#

require File.expand_path(File.dirname(__FILE__) + '/../config')
require File.expand_path(File.dirname(__FILE__) + '/book.inc')

# get an instance of the SDB interface using the default configuration
sdb = AWS::SimpleDB.new :region => 'us-east-1'

Dir.entries('.').select{ |file| file.match(/^[a-zA-Z0-9_-]*\.rb$/)}.each do |filename|
  File.open filename do |f|
    data = f.read
    size = f.size
    md5 = Digest::MD5.hexdigest(data)

    item = sdb.domains[BOOK_FILE_DOMAIN].items[filename]
    item.attributes.put(
      add: { name: filename }, # Append values to this attribute
      replace: { md5: md5, size: "%08d" % size}  # Set all values for this attribute, replacing current values
    )
  end
end
\end{lstlisting}
\end{itemize}
\end{frame}

%%%%%%%%%%%%%%%%%%%%%%%%%%%%%%%%%%%%%%%%%%%%%%%%%%%%%%%%%%%%%%%%%%%%%%%%%%%%%%

%%%%%%%%%%%%%%%%%%%%%%%%%%%%%%%%%%%%%%%%%%%%%%%%%%%%%%%%%%%%%%%%%%%%%%%%%%%%%%
\begin{frame}[fragile, allowframebreaks]
\frametitle{Storing Multiple Items Efficiently}
\begin{itemize}
\item Another SimpleDB function, \texttt{batch\_put\_attributes}, can be used to perform multiple
\texttt{put} attributes calls with a single request. This can be faster and more efficient
than making a long series of individual requests. Here's a modified version of the
previous example that uses the batch requests

\lstset{language=Ruby, style=eclipse}
\begin{lstlisting}
#
# aws_app/simpledb/batch_insert_items.rb
#
# Usage: ruby simpledb/batch_insert_items.rb
#

require File.expand_path(File.dirname(__FILE__) + '/../config')
require File.expand_path(File.dirname(__FILE__) + '/book.inc')

# get an instance of the SDB interface using the default configuration
sdb = AWS::SimpleDB.new :region => 'us-east-1'

items = []

Dir.entries('.').select{ |file| file.match(/^[a-zA-Z0-9_-]*\.rb$/)}.each do |filename|
  File.open filename do |f|
    data = f.read
    size = f.size
    md5 = Digest::MD5.hexdigest data

    items << { name: filename,
               attributes: [ {name: 'name', value: filename, replace: false},
                             {name: 'md5', value: md5, replace: true},
                             {name: 'size', value: "%08d" % size, replace: true}
               ]
    }
  end
  if items.size == 10
    response = sdb.client.batch_put_attributes domain_name: BOOK_FILE_DOMAIN, items: items
    puts response.data.to_s
    items.clear
  end
end

if items.size > 0
  response = sdb.client.batch_put_attributes domain_name: BOOK_FILE_DOMAIN, items: items
  puts response.data.to_s
end
\end{lstlisting}
\end{itemize}
\end{frame}
%%%%%%%%%%%%%%%%%%%%%%%%%%%%%%%%%%%%%%%%%%%%%%%%%%%%%%%%%%%%%%%%%%%%%%%%%%%%%%

%%%%%%%%%%%%%%%%%%%%%%%%%%%%%%%%%%%%%%%%%%%%%%%%%%%%%%%%%%%%%%%%%%%%%%%%%%%%%%
\begin{frame}[fragile, allowframebreaks]
\frametitle{Running a Query}
\begin{itemize}
\item With the data stored, running queries is simple:

\lstset{language=Ruby, style=eclipse}
\begin{lstlisting}
#
# aws_app/simpledb/query_domain.rb
#
# Usage: ruby simpledb/query_domain.rb
#

require File.expand_path(File.dirname(__FILE__) + '/../config')
require File.expand_path(File.dirname(__FILE__) + '/book.inc')

# get an instance of the SDB interface using the default configuration
sdb = AWS::SimpleDB.new :region => 'us-east-1'

# seleccionar los archivos cuyo nombre empiece por list
# query = "select * from #{BOOK_FILE_DOMAIN} where name like 'list%'"

# to select files with size > 1000 bytes
query = "select * from #{BOOK_FILE_DOMAIN} where size > '00001000'"

begin
  result = sdb.client.select select_expression: query
  puts result.data[:items]
end while result.data[:next_token]
\end{lstlisting}
\end{itemize}
\end{frame}

%%%%%%%%%%%%%%%%%%%%%%%%%%%%%%%%%%%%%%%%%%%%%%%%%%%%%%%%%%%%%%%%%%%%%%%%%%%%%%
%%%%%%%%%%%%%%%%%%%%%%%%%%%%%%%%%%%%%%%%%%%%%%%%%%%%%%%%%%%%%%%%%%%%%%%%%%%%%%
\begin{frame}[fragile, allowframebreaks]
\frametitle{Advanced Queries}
\begin{itemize}
\item SimpleDB queries have the following general form:
\lstset{language=Ruby}
\begin{lstlisting}
select output_list
from domain_name
[where expresion]
[order by clause]
[limit clause]
\end{lstlisting}

\item The optional \texttt{where} expressions are a combination of \texttt{attribute names}, \texttt{constant
values}, \texttt{item name}, and \texttt{operators}. The following list shows the most used operators that
can be used in the \texttt{where} expression:
\begin{itemize}
\item \texttt{=, !=, >, >=, <, <=, and, or}
\item \texttt{like} (contains the specified constant, with “\%” as a wildcard,) \texttt{not like}
\item \texttt{between/and} (In a specified range), \texttt{in} (Matches an item from a list)
\item \texttt{is null } (Attribute does not exist), \texttt{is not null}, \texttt{every} (Every value of a multi-valued attribute satisfies the expression)
\end{itemize}

\item We would like to execute 
\lstset{language=shell}
\begin{lstlisting}
$ ruby simpledb/query_domain_cmd.rb "size > '00001000'"
\end{lstlisting}

\lstset{language=Ruby, style=eclipse}
\begin{lstlisting}
#
# aws_app/simpledb/query_domain_cmd.rb
#
# Usage: ruby simpledb/query_domain_cmd.rb "size > '00001000'"
#

require File.expand_path(File.dirname(__FILE__) + '/../config')
require File.expand_path(File.dirname(__FILE__) + '/book.inc')

# get an instance of the SDB interface using the default configuration
sdb = AWS::SimpleDB.new :region => 'us-east-1'

query = "select * from #{BOOK_FILE_DOMAIN.to_s} "
query += "where #{ARGV.join ' '}" unless ARGV.size == 0

puts "Final query: #{query}"

begin
  result = sdb.client.select select_expression: query
  puts result.data[:items]
end while result.data[:next_token]
\end{lstlisting}
\end{itemize}
\end{frame}

%%%%%%%%%%%%%%%%%%%%%%%%%%%%%%%%%%%%%%%%%%%%%%%%%%%%%%%%%%%%%%%%%%%%%%%%%%%%%%
%%%%%%%%%%%%%%%%%%%%%%%%%%%%%%%%%%%%%%%%%%%%%%%%%%%%%%%%%%%%%%%%%%%%%%%%%%%%%%
\begin{frame}[fragile, allowframebreaks]
\frametitle{Augmenting Items with Additional Data}
\begin{itemize}
\item Once an item has been created, additional calls to the \texttt{put\_attributes} method can
be used to create additional attributes or additional values for an existing attribute.
Let’ s augment each item in the \texttt{BOOK\_FILE\_DOMAIN with the file’ s modification time}
using this script:

\lstset{language=Ruby, style=eclipse}
\begin{lstlisting}
#
# aws_app/simpledb/augmented_items.rb
#
# Usage: ruby simpledb/augmented_items.rb
#

require File.expand_path(File.dirname(__FILE__) + '/../config')
require File.expand_path(File.dirname(__FILE__) + '/book.inc')

# get an instance of the SDB interface using the default configuration
sdb = AWS::SimpleDB.new :region => 'us-east-1'

query = "select name from #{BOOK_FILE_DOMAIN}"

begin
  result = sdb.client.select(select_expression: query)
  result.data[:items].each do |item|
    mod_time = File.mtime item[:name]
    puts item[:name], mod_time.to_s

    db_item = sdb.domains[BOOK_FILE_DOMAIN].items[item[:name]]
    db_item.attributes.put( replace: { mod_time: mod_time.to_s } )
  end
end while result.data[:next_token]

\end{lstlisting}
\end{itemize}
\end{frame}

%%%%%%%%%%%%%%%%%%%%%%%%%%%%%%%%%%%%%%%%%%%%%%%%%%%%%%%%%%%%%%%%%%%%%%%%%%%%%%
%%%%%%%%%%%%%%%%%%%%%%%%%%%%%%%%%%%%%%%%%%%%%%%%%%%%%%%%%%%%%%%%%%%%%%%%%%%%%%
\begin{frame}[fragile, allowframebreaks]
\frametitle{Augmenting Items with Additional Data}
\begin{itemize}
\item Once an item has been created, additional calls to the \texttt{put\_attributes} method can
be used to create additional attributes or additional values for an existing attribute.
Let’ s augment each item in the \texttt{BOOK\_FILE\_DOMAIN with the file’ s modification time}
using this script:

\lstset{language=Ruby, style=eclipse}
\begin{lstlisting}
#
# aws_app/simpledb/augmented_items.rb
#
# Usage: ruby simpledb/augmented_items.rb
#

require File.expand_path(File.dirname(__FILE__) + '/../config')
require File.expand_path(File.dirname(__FILE__) + '/book.inc')

# get an instance of the SDB interface using the default configuration
sdb = AWS::SimpleDB.new :region => 'us-east-1'

query = "select name from #{BOOK_FILE_DOMAIN}"

begin
  result = sdb.client.select(select_expression: query)
  result.data[:items].each do |item|
    mod_time = File.mtime item[:name]
    puts item[:name], mod_time.to_s

    db_item = sdb.domains[BOOK_FILE_DOMAIN].items[item[:name]]
    db_item.attributes.put( replace: { mod_time: mod_time.to_s } )
  end
end while result.data[:next_token]

\end{lstlisting}
\end{itemize}
\end{frame}

%%%%%%%%%%%%%%%%%%%%%%%%%%%%%%%%%%%%%%%%%%%%%%%%%%%%%%%%%%%%%%%%%%%%%%%%%%%%%%
%%%%%%%%%%%%%%%%%%%%%%%%%%%%%%%%%%%%%%%%%%%%%%%%%%%%%%%%%%%%%%%%%%%%%%%%%%%%%%
\begin{frame}[fragile, allowframebreaks]
\frametitle{Storing Multiple Values for an Attribute}
\begin{itemize}
\item Storing multiple values for a single attribute is very easy; there are two ways to do
this.
\begin{itemize}
\item You can store multiple modification times using the code from the last section, with repeated
calls like so

\lstset{language=Ruby, style=eclipse}
\begin{lstlisting}[escapechar=!]
#
# aws_app/simpledb/augmented_items.rb
#
# Usage: ruby simpledb/augmented_items.rb
#
!\vdots!
    # A hash of attribute names and values to append to this item
    db_item.attributes.put( add: { mod_time: mod_time.to_s } )
!\vdots!
\end{lstlisting}

Then update modifed time and run the script

\lstset{language=shell}
\begin{lstlisting}
$ touch *.rb
$ ruby simpledb/augmented_items.rb
\end{lstlisting}

\item You can store sever al values at the same time like this:
\begin{lstlisting}[escapechar=!]
#
# aws_app/simpledb/augmented_items.rb
#
# Usage: ruby simpledb/augmented_items.rb
#
!\vdots!
    db_item.attributes.put( replace: { mod_time: [mod_time.to_s, Time.now] } ) !\circle{1}!
!\vdots!
\end{lstlisting}


\end{itemize}
\end{itemize}
\end{frame}

%%%%%%%%%%%%%%%%%%%%%%%%%%%%%%%%%%%%%%%%%%%%%%%%%%%%%%%%%%%%%%%%%%%%%%%%%%%%%%
%%%%%%%%%%%%%%%%%%%%%%%%%%%%%%%%%%%%%%%%%%%%%%%%%%%%%%%%%%%%%%%%%%%%%%%%%%%%%%
\begin{frame}[fragile, allowframebreaks]
\frametitle{Deleting Attributes}
\begin{itemize}
\item If you’ve no more need for an attribute, you can delete attributes it like this:

\lstset{language=Ruby, style=eclipse}
\begin{lstlisting}[escapechar=!]
#
# aws_app/simpledb/delete_attributes.rb
#
# Usage: ruby simpledb/delete_attributes.rb
#

require File.expand_path(File.dirname(__FILE__) + '/../config')
require File.expand_path(File.dirname(__FILE__) + '/book.inc')

# get an instance of the SDB interface using the default configuration
sdb = AWS::SimpleDB.new :region => 'us-east-1'

query = "select name from #{BOOK_FILE_DOMAIN}"

attributes = %W{mod_time flavor}
begin
  result = sdb.client.select(select_expression: query)
  result.data[:items].each do |item|

    db_item = sdb.domains[BOOK_FILE_DOMAIN].items[item[:name]]
    db_item.attributes.delete(attributes )
  end
end while result.data[:next_token]
\end{lstlisting}
\end{itemize}

\end{frame}

%%%%%%%%%%%%%%%%%%%%%%%%%%%%%%%%%%%%%%%%%%%%%%%%%%%%%%%%%%%%%%%%%%%%%%%%%%%%%%
%%%%%%%%%%%%%%%%%%%%%%%%%%%%%%%%%%%%%%%%%%%%%%%%%%%%%%%%%%%%%%%%%%%%%%%%%%%%%%
\begin{frame}[fragile, allowframebreaks]
\frametitle{Deleting Items}
\begin{itemize}
\item You can also delete an entire item, by deleting all the attributes. Here’ s how to do it:

\lstset{language=Ruby, style=eclipse}
\begin{lstlisting}[escapechar=!]
#
# aws_app/simpledb/delete_items.rb
#
# Usage: ruby simpledb/delete_items.rb
#

require File.expand_path(File.dirname(__FILE__) + '/../config')
require File.expand_path(File.dirname(__FILE__) + '/book.inc')

# get an instance of the SDB interface using the default configuration
sdb = AWS::SimpleDB.new :region => 'us-east-1'

query = "select * from #{BOOK_FILE_DOMAIN} where name = 'config.rb'"

begin
  result = sdb.client.select(select_expression: query)
  result.data[:items].each do |item|
    db_item = sdb.domains[BOOK_FILE_DOMAIN].items[item[:name]]
    db_item.attributes.delete( db_item.attributes.to_h )
  end
end while result.data[:next_token]
\end{lstlisting}
\end{itemize}

\end{frame}

%%%%%%%%%%%%%%%%%%%%%%%%%%%%%%%%%%%%%%%%%%%%%%%%%%%%%%%%%%%%%%%%%%%%%%%%%%%%%%
%%%%%%%%%%%%%%%%%%%%%%%%%%%%%%%%%%%%%%%%%%%%%%%%%%%%%%%%%%%%%%%%%%%%%%%%%%%%%%
\begin{frame}[fragile, allowframebreaks]
\frametitle{Monitoring Domain Statistics}
\begin{itemize}
\item Here’ s a script that displays metadata domains:

\lstset{language=Ruby, style=eclipse}
\begin{lstlisting}
#
# aws_app/simpledb/metadata.rb
#
# Usage: ruby simpledb/metadata.rb
#

require File.expand_path(File.dirname(__FILE__) + '/../config')
require File.expand_path(File.dirname(__FILE__) + '/book.inc')

# get an instance of the SDB interface using the default configuration
sdb = AWS::SimpleDB.new :region => 'us-east-1'

sdb.domains.each do |domain|
  puts "Domain: #{domain.name}: #{domain.metadata.to_h.to_s}"
end
\end{lstlisting}


\lstset{language=shell}
\begin{lstlisting}[escapechar=!]
$ ruby simpledb/metadata.rb
!\outputcommand{Domain: feed\_items: {:item\_count=>0, :item\_names\_size\_bytes=>0, :attribute\_name\_count=>0, :attribute\_names\_size\_bytes=>0, :attribute\_value\_count=>0, :attribute\_values\_size\_bytes=>0, :timestamp=>1394701235}\\
Domain: feeds: {:item\_count=>0, :item\_names\_size\_bytes=>0, :attribute\_name\_count=>0, :attribute\_names\_size\_bytes=>0, :attribute\_value\_count=>0, :attribute\_values\_size\_bytes=>0, :timestamp=>1394701235}\\
Domain: files: {:item\_count=>1, :item\_names\_size\_bytes=>9, :attribute\_name\_count=>3, :attribute\_names\_size\_bytes=>11, :attribute\_value\_count=>3, :attribute\_values\_size\_bytes=>49, :timestamp=>1394701235}}
\end{lstlisting}

\end{itemize}

\end{frame}

%%%%%%%%%%%%%%%%%%%%%%%%%%%%%%%%%%%%%%%%%%%%%%%%%%%%%%%%%%%%%%%%%%%%%%%%%%%%%%
\section{Homework}
\begin{frame}[fragile, allowframebreaks]
\frametitle{Homework}

\begin{itemize}
\item Dado el localizador url de un sitio registrar la frecuencia de uso de todas las palabaras utilizadas como títulos en todas las páginas internas accesibles desde ese localizador. El título en un documento html está escrito entre etiqutas \verb%<h1></h1>% o \verb%<h2></h2>%, etc. Una posible diseño del servicio sería el siguiente:

\begin{enumerate}
\item 4 colas: url\_queue, title\_queue, url\_crawler\_queue, counter\_queue
\item Un primer script que simplemente añade la url de partida (por ejemplo htttp://www.elpais.es) a la cola url\_queue
\item Un proceso que recoge los mensajes en la cola url\_queue y descarga la página html al S3 y genera dos mensajes con la información de la localización
de la página: una para la cola title\_queue y otro para la cola  url\_crawler\_queue
\begin{enumerate}
\item Otro proceso que recoge los mensajes en la cola title\_queue  y crea otro fichero S3 con texto plano con el contenido de los títulos y añade un mensaje a la cola counter\_queue
\item Otro proceso que paralelamente al anterior coge url\_crawler\_queue busca url relativas al sitio y las añade a la cola url\_queue
\end{enumerate}
\item Un proceso más que toma el texto plano del fichero en el S3 indicado en el mensaje del la cola counter\_queue y registar en simpleDB la utilización de las palabras
\end{enumerate}

\item Detalles a tener en cuenta:
\begin{itemize}
\item Una página descargada debe marcarse con un item en simpleDB como tal para evitar bucles de descargas infinitos.
\item ¿Transacciones en simpleDB? Cuando del último proceso hayan más de una instancia corriendo una implementación sin transacciones no funcionará. Ver Conditional Puts/Deletes
\end{itemize}
\end{itemize}
\end{frame}
%%%%%%%%%%%%%%%%%%%%%%%%%%%%%%%%%%%%%%%%%%%%%%%%%%%%%%%%%%%%%%%%%%%%%%%%%%%%%%
\end{document}
